\documentclass[11pt]{article}
\usepackage{amsmath}
\usepackage{amsfonts}
\usepackage{amsthm}
\usepackage[utf8]{inputenc}
\usepackage[margin=0.75in]{geometry}

\title{CSC111 Winter 2026 Project 1}
\author{Alsade Brianna Daley, Elif Cakici, Zila Okoroze-Trust}
\date{\today}

\begin{document}
\maketitle

\section{Running the game}
The game can be run by executing \texttt{adventure.py} using Python, no additional libraries required.

\section*{Game Map}
Example game map (replace with actual IDs from \texttt{game\_data.json}):

\begin{verbatim}
 -1  1  2
 -1  3 -1
  5  4  7  9 10
 -1  6  8 -1 11
\end{verbatim}

Starting location is: 1

\section*{Game solution}
List of commands to win the game:
\begin{enumerate}
    \item pick up toonie
    \item pick up pencil
    \item go east
    \item pick up usb drive
    \item go west
    \item go south
    \item pick up jam
    \item go south
    \item go west
    \item pick up key fob  : answer the riddle with "robarts"
    \item go east
    \item go south
    \item use jam
    \item pick up lucky mug
    \item go north
    \item go east
    \item go east
    \item use pencil
    \item pick up laptop charger
    \item go east
    \item go south
    \item use key fob
\end{enumerate}

\end{lstlisting}

\section*{Lose condition(s)}
Description: The player loses the game if they exceed the maximum allowed number of moves, which is 35.

List of commands that cause a loss are: (example from \texttt{lose\_demo}):

\begin{verbatim}
["go south", "go north"] * 17 + ["go south"]
\end{verbatim}

Parts of the code involved:
\begin{itemize}
    \item \texttt{adventure.py}, \texttt{AdventureGame} class: the \texttt{player.moves} attribute tracks the number of moves.
    \item The game loop checks if moves exceed \texttt{MAX\_MOVES} and if so, ends the game.
    \item Helper functions in \texttt{adventure.py} assist with validating moves and commands.
\end{itemize}

\section*{Inventory}

\begin{enumerate}
    \item All location IDs that involve items in the game: 1, 2, 3, 5, 6, 8, 9
    \item Item data:
    \begin{enumerate}
        \item Item 1:
        \begin{itemize}
            \item Item name: "toonie"
            \item Item start location ID: 1
            \item Item target location ID: 4
        \end{itemize}
        \item Item 2:
        \begin{itemize}
            \item Item name: "pencil"
            \item Item start location ID: 1
            \item Item target location ID: 9
        \end{itemize}
        \item Item 3:
        \begin{itemize}
            \item Item name: "usb drive"
            \item Item start location ID: 2
            \item Item target location ID: 11
        \end{itemize}
        \item Item 4:
        \begin{itemize}
            \item Item name: "jam"
            \item Item start location ID: 3
            \item Item target location ID: 6
        \end{itemize}
        \item Item 5:
        \begin{itemize}
            \item Item name: "key fob"
            \item Item start location ID: 5
            \item Item target location ID: 11
        \end{itemize}
        \item Item 6:
        \begin{itemize}
            \item Item name: "lucky mug"
            \item Item start location ID: 6
            \item Item target location ID: 11
        \end{itemize}
        \item Item 7:
        \begin{itemize}
            \item Item name: "laptop charger"
            \item Item start location ID: 9
            \item Item target location ID: 11
        \end{itemize}
    \end{enumerate}

    \item Example commands to pick up and use items (from \texttt{inventory\_demo}):
    \begin{verbatim}
"pick up toonie",
"pick up pencil",
"inventory",
"go east",
"pick up usb drive",
"inventory"
"drop pencil"
    \end{verbatim}

    \item Code involved in handling the \texttt{inventory} command:
    \begin{itemize}
        \item \texttt{adventure.py}, the function \texttt{handle\_menu\_command()}
        \item \texttt{Player} class in \texttt{game\_entities.py}, methods \texttt{pick\_up}, \texttt{drop\_item}, \texttt{has\_item} 
        \item Helper functions in \texttt{adventure.py} assist with validating item commands
    \end{itemize}
\end{enumerate}

\section*{Score}

\begin{enumerate}
    \item Players earn points by depositing items in their target locations. The first location to earn points is location 4. Commands leading to score increase (from \texttt{scores\_demo}):
\begin{verbatim}
"pick up toonie",
"go south",
"score"
\end{verbatim}

    \item Code involved:
    \begin{itemize}
        \item \texttt{adventure.py}, \texttt{handle\_item\_command()}, updates \texttt{player.score}
        \item \texttt{handle\_menu\_command()} displays the current score
        \item \texttt{Player} class in \texttt{game\_entities.py}
        \item Helper functions in \texttt{adventure.py} are used for validating item locations and scoring
    \end{itemize}
\end{enumerate}

\section*{Enhancements}
\begin{enumerate}
    \item Player Class Puzzle Enhancement
    \begin{itemize}
        \item Brief description: Added a system for using items to solve interconnected puzzles in the game. Steps include:
        \begin{enumerate}
            \item Give the beggar a \texttt{toonie} to receive a hint
            \item Solve a riddle to obtain the \texttt{key fob}
            \item Exchange \texttt{jam} or cookies for a \texttt{lucky mug} in the dining hall
            \item Be tricked into joining a cookie line, losing extra steps
            \item Give a pencil to the receptionist to get a charger in fewer steps
            \item Use the \texttt{key fob} to unlock the dorm door
        \end{enumerate}
        \item Complexity level: Medium
        \item Justification: Requires significant additions to the baseline game logic, including:
        \begin{itemize}
            \item Updating the \texttt{Player} class to handle \texttt{use item} functionality
            \item Adding multiple interconnected puzzle steps
            \item Modifying \texttt{adventure.py} for item interactions
            \item Adding new items
            \item Adding code to manage various types of wrong inputs
            \item Managing game state, unlocked locations, and move count
            \item Challenge: Attempted to move helper functions to a separate file, but circular import errors occurred, so logic remains in \texttt{adventure.py}
        \end{itemize}
        \item Code involved:
        \begin{itemize}
            \item \texttt{adventure.py}, \texttt{AdventureGame} class, methods handling item interactions and location puzzles
            \item \texttt{Player} class in \texttt{game\_entities.py} for inventory and item usage, including interact modules used to manage puzzle locations
            \item \texttt{handle\_command()} and \texttt{handle\_item\_command()} for executing the \texttt{use item} commands
            \item Helper functions in \texttt{adventure.py} used to validate moves, items, and puzzle conditions
        \end{itemize}
        \item Commands from \texttt{enhancements\_demo}:
\begin{verbatim}
"pick up toonie",
"go south",
"go south",
"go west",
"pick up key fob"
\end{verbatim}
    \end{itemize}
    \end{itemize}

\end{document}
